%------------------------------------------------------------------------------
% \PART 6
%------------------------------------------------------------------------------
\chapter[Анализ результатов тестирования IP-блока нейронной сети для распознавания рукописных цифр]
{АНАЛИЗ РЕЗУЛЬТАТОВ ТЕСТИРОВАНИЯ IP-БЛОКА НЕЙРОННОЙ СЕТИ ДЛЯ РАСПОЗНАВАНИЯ РУКОПИСНЫХ ЦИФР}

%------------------------------------------------------------------------------
% \PART 6.1 Text
%------------------------------------------------------------------------------
\section{Описание среды тестирования}\par
\hspace*{12.5 mm}В качестве аппаратной платформы для реализации и тестирования 
разработанных решений использовалась плата на базе программируемой логической 
интегральной схемы ZYBO Z7 от компании Digilent. Плата ZYBO Z7 представляет 
собой одноплатный компьютер, основанный на СнК Xilinx Zynq-7000, которая 
сочетает в себе двухъядерный процессор ARM Cortex-A9 и среду программируемой 
логики (PL) на одном чипе\cite{ZYBO}.

ПЛИС ZYBO Z7 оснащена широким набором периферийных интерфейсов, таких как HDMI, 
USB, Ethernet, аудиоразъёмы и слот для карты памяти microSD. Такая конфигурация
позволяет использовать плату как для высокопроизводительных вычислений, так и 
для задач обработки сигналов и изображений в реальном времени. Программируемая 
часть ПЛИС используется для создания специализированных аппаратных ускорителей 
обработки данных, а процессорная часть (Processing System, PS) обеспечивает 
управление и взаимодействие между различными компонентами системы\cite{ZYNQ}.

Для упрощения разработки приложений и взаимодействия между процессорной 
системой и пользовательскими IP-блоками программируемой логики, в работе 
применялась среда PYNQ (Python Productivity for Zynq). 

Общая структура использования PYNQ для разработки Zynq представлена на 
рисунке~\ref{fig:pynq-struct} \cite{PYNQ-IMG}.

\insertfigure{pynq-struct}{pic/pynq.png}{Зависимость точности от разрядности}{12.5cm}

PYNQ представляет собой 
специализированную программную платформу, основанную на операционной системе 
Linux, которая позволяет использовать язык программирования Python для 
управления аппаратными ресурсами Zynq\cite{PYNQ}. PYNQ предоставляет удобный интерфейс
для работы с IP-ядрами, созданными в среде Vivado, через использование 
библиотек и драйверов высокого уровня.

Программная часть проекта выполнялась на процессоре ARM, работающем под 
управлением операционной системы PYNQ Linux. Управление аппаратными IP-блоками 
осуществлялось через Python API, что позволило оперативно тестировать и 
модифицировать аппаратные модули без необходимости перекомпиляции всего 
проекта. Благодаря этому подходу существенно ускорился процесс отладки и 
тестирования системы.

%------------------------------------------------------------------------------
% \PART 6.2 Text
%------------------------------------------------------------------------------
\section{Тестирование разработанного IP-блока нейронной сети}\par
\hspace*{12.5 mm}Тестирование осуществлялось посредством запуска специально 
разработанного скрипта на языке Python в среде PYNQ.

Основной задачей тестирования являлась проверка корректности выполнения 
операций нейронной сети, реализованных в аппаратной логике, и 
взаимодействия IP-блока с процессорной системой через интерфейсы AXI4-lite и uP. 
Для этого использовался скрипт на Python, который выполнял следующие функции: 

    1 Инициализация IP-блока нейронной сети. 

    2 Загрузка тестовых данных во входные регистры IP-блока. 
    
    3 Запуск аппаратной обработки данных. 
    
    4 Считывание результатов вычислений из выходных регистров. 
    
Также отдельно осуществляется сравнение полученных результатов с эталонными 
значениями для верификации корректности работы. 

На этапе инициализации скрипт с помощью предоставленных библиотек PYNQ, таких 
как {Overlay}, загружает конфигурацию программируемой логики и получал доступ 
к IP-блоку.

Для подачи входных данных использовался архив с тестовыми изображениями из 
базы MNIST. Осущесмтвляется чтение изображения и значения. Затем изображение
нормализуется в соответствии с описанными ранее треброваниями. Скрипт 
производит запись данных во входные регистры через AXI-интерфейс и
инициировал начало вычислений путём установки управляющих флагов.

По завершении обработки скрипт ожидает сигнал готовности от IP-блока. После 
этого результаты обработки считывались из выходных регистров. Полученные данные 
сравнивались со считанным значением из базы.

В тестовом окружении необходимо подключить библиотеку PYNQ и файл с прошивкой:

\small
\fontsize{12pt}{12pt}\selectfont
\begin{verbatim}
    from pynq import Overlay 

    platform = Overlay("design_1.bit")
    neural_net = platform.m_nn_0
\end{verbatim}
\normalsize
\fontsize{14pt}{14pt}\selectfont

Полный скрипт по обработке изображения на ПЛИС представлен в приложении Д.

На рисунке~\ref{fig:hw-res} представлен пример работы данного скрипта.

\insertfigurescustom{pic/ex2_fpga.jpg}{pic/ex3_fpga.jpg}{6.5cm}{Результаты распознования рукописных цифр}{hw-res}{sidebyside}

%------------------------------------------------------------------------------
% \PART 6.3 Text
%------------------------------------------------------------------------------
\section{Экспериментальное исследование IP-блока нейронной сети}\par
\hspace*{12.5 mm}Для оценки эффективности работы разработанного IP-блока 
нейронной сети было проведено три последовательных эксперимента. Эксперименты 
направлены на проверку корректности работы системы, оценку точности 
классификации и исследование влияния разрядности весовых коэффициентов на 
качество работы сети. Произведено три эксперимента:

    1 Тестирование на 10000 изображений из MNIST. Построение матрицы спутывания
    и сравнение с эталоном.

    2 Использование консольного приложения по рисованию на холсте с 
    возможностью сохранить изображение в формате 28 на 28 и дальнейшей 
    обработке на ПЛИС.

    3 Изменение разрядности представления весовых коэффициентов и проверка 
    зависимости точности от разрядности.

    Скрипт представлен в приложении Д.

%------------------------------------------------------------------------------
% \PART 6.3.1 Text
%------------------------------------------------------------------------------
\anonsection{}Первым этапом тестирования стала проверка IP-блока на большом 
количестве стандартных данных. Для этого использовалась тестовая выборка из 
10000 изображений базы данных MNIST, содержащей изображения рукописных цифр 
размером $28\times28$ пикселей.

Каждое изображение подавалось на вход разработанному IP-блоку, после чего 
считывался результат классификации. На основании полученных ответов была 
построена матрица спутывания, отражающая количество верных и ошибочных 
классификаций для каждого класса.

Матрица спутывания позволяет оценить, какие именно классы наиболее часто 
путаются между собой, а также вычислить стандартные метрики качества 
классификации, такие как точность, полнота и 
специфичность.

Сравнение результатов работы IP-блока с эталонной программной моделью нейронной 
сети показало полное соответствие, что свидетельствует о правильности 
аппаратной реализации.

На рисунке~\ref{fig:conf-m} представлены матрицы спутывания.

\insertfigurescustom{pic/conf-matrix.jpg}{pic/conf-matrix.jpg}{8cm}{Аппаратная и эталонная матрицы спутывания}{conf-m}{sidebyside}

Если проанализировать результаты классификации при реализации модели в 
фиксированной точности и с плавающей запятой, можно отметить, что точность 
распознавания снизилась с 98.02\% до 97.9\%. Это соответствует потере базовой 
точности на 0.12\%.

Данное снижение связано с квантованием весовых коэффициентов, которое приводит 
к небольшому искажению параметров модели. Однако потеря точности оказалась 
незначительной, что свидетельствует о высокой устойчивости предложенной 
архитектуры к операциям квантования.

Результаты представлены в графическом материале ГУИР 431289.007 ПЛ.

%------------------------------------------------------------------------------
% \PART 6.3.2 Text
%------------------------------------------------------------------------------
\anonsection{}Вторым этапом исследования стало тестирование на изображениях 
цифр, созданных вручную в специально разработанном консольном приложении для 
рисования. Программа позволяла пользователю нарисовать цифру на виртуальном 
холсте, сохранить изображение в формате $28\times28$ пикселей.

Затем отдельно запускается скрипт по обработке изображения на ПЛИС.

Данный эксперимент позволил оценить устойчивость работы IP-блока к вариациям 
входных данных, отличающихся от стандартных изображений MNIST. 

Отрисовано 100 различных изображений и проверена работа данного IP-блока с 
построением матрицы спутывания и вычислением точности.

На рисунке~\ref{fig:conf-m1} представлена матрица спутывания.

\insertfigure{conf-m1}{pic/conf-matrix.jpg}{Матрица спутывания на не стандартных данных}{10cm}

Результаты представлены в графическом материале ГУИР 431289.007 ПЛ.

%------------------------------------------------------------------------------
% \PART 6.3.3 Text
%------------------------------------------------------------------------------
\anonsection{}Заключительный эксперимент был направлен на изучение влияния 
разрядности представления весовых коэффициентов на качество классификации. Для 
этого были проведены тестирования с различными значениями битности 
представления весов, такими как 16 бит, 12 бит, 8 бит и т.д.

Для повышения модульности и удобства в разработке 
использовался пользовательский пакет, содержащий локальные параметры, такие как 
разрядность данных. Это позволило централизованно управлять настройками всех 
модулей нейросетевой модели и обеспечило их согласованность между собой. 
Использование пакета упростило процесс масштабирования и модификации 
архитектуры, так как изменения параметров выполняются в одном месте без 
необходимости правки каждого отдельного модуля.

\small
\fontsize{12pt}{12pt}\selectfont
\begin{verbatim}
package nn_param_pkg;

    localparam INT_BITS  = 6;
    localparam FRC_BITS  = 7;
    localparam BITS      = 13;
    
    localparam PICT_SIZE = 28;
    localparam WIDTH     = 784;
  
    localparam ADDR_RC   = 5;
    localparam ADDR_OUT  = 10;
  
    localparam RCO_TYPE  = 2;
  
    localparam NN_FSM_BUS_WIDTH  = 5;
      
endpackage
\end{verbatim}
\normalsize
\fontsize{14pt}{14pt}\selectfont

  На рисунке~\ref{fig:acc} представлен график зависимости точности от 
разрядности.

\insertfigure{acc}{pic/Acc_LUTs_FFs.jpg}{Зависимость точности от разрядности}{16cm}

На каждом этапе тестирования сети применялись квантизованные веса с 
соответствующей разрядностью. Далее выполнялось тестирование на стандартной 
выборке MNIST, с построением матриц спутывания и вычислением точности 
классификации.

Анализ результатов показал, что при снижении разрядности до 8 бит потери 
точности классификации остаются незначительными, что подтверждает эффективность 
использования фиксированной арифметики для оптимизации аппаратной реализации. 
Однако дальнейшее уменьшение разрядности до 6 бит приводило к заметному 
ухудшению качества распознавания.

Результаты представлены в графическом материале ГУИР 431289.007 ПЛ.

%------------------------------------------------------------------------------
% \PART 6.4 Text
%------------------------------------------------------------------------------
\section{Сравнение с аналогичнми архитектурами}\par
\hspace*{12.5 mm}Для оценки эффективности предложенного IP-блока нейронной сети 
на базе слоя LST было проведено сравнение с рядом существующих решений, 
предназначенных для задачи классификации изображений рукописных цифр из базы
MNIST.

\begin{table*}[ht]
    Таблица 4 – Сравнение архитектур нейронных сетей для задачи классификации \\\hspace*{2.5cm}рукописных цифр из базы MNIST\par
    \renewcommand{\arraystretch}{1.2}
    \begin{tabularx}{\textwidth}{|>{\raggedright\arraybackslash}X|>{\centering\arraybackslash}X|>{\centering\arraybackslash}X|>{\centering\arraybackslash}X|}
      \hline
      {Автор} & {Архитектура нейронной сети} & {Число параметров} & {Точность} \\ \hline
      Liang, et al. (2018)\cite{Liang} & 784-2048-2048-2048-10 & 10 100 000 & 98.32\% \\ \hline
      Umuroglu, et al. (2017)\cite{Umuroglu} & 784-1024-1024-10 & 1 863 690 & 98.40\% \\ \hline
      Medus, et al. (2019)\cite{Medus} & 784-600-600-10 & 891 610 & 98.63\% \\ \hline
      Huynh\cite{Huynh} & 784-126-126-10 & 115 920 & 98.16\% \\ \hline
      Huynh\cite{Huynh} & 784-40-40-40-10 & 34 960 & 97.20\% \\ \hline
      Westby, et al.\cite{Westby} & 784-12-10 & 9 550 & 93.25\% \\ \hline
      LST-1 & LST-784-10 & 9474 & 97.9\% \\ \hline
    \end{tabularx}
\end{table*}

Из таблицы видно, что большинство существующих архитектур нейронных сетей, 
ориентированных на высокую точность распознавания MNIST, характеризуются 
большим числом обучаемых параметров, что приводит к увеличению затрат ресурсов 
при аппаратной реализации. Например, нейронная сеть Liang et al.\cite{Liang} 
имеет более 10 миллионов параметров, а сеть Umuroglu et al.\cite{Umuroglu} – 
порядка 1.8 миллиона параметров.

Предложенная архитектура LST-1 отличается принципиально меньшим числом 
параметров – всего 9474, что делает её крайне привлекательной для реализации в 
ресурсно-ограниченных средах, таких как ПЛИС или встроенные системы. При этом 
достигается высокая точность распознавания — 97.9\%, что сравнимо с более 
крупными архитектурами.

Таким образом, аппаратная реализация нейронной сети на базе слоя LST 
демонстрирует компромисс между компактностью модели и точностью классификации. 
Это позволяет использовать её в задачах, где критичны малое энергопотребление, 
быстродействие и ограниченный объем ресурсов.

Предложенный IP-блок LST может рассматриваться как альтернатива классическим 
полносвязным слоям при разработке нейронных сетей прямого распространения 
(многослойных перцептронов, MLP), особенно в случае, если требуется минимизация 
площади кристалла или энергопотребления в встраиваемых устройствах.

